\documentclass[a4paper,11pt]{article}

\usepackage[utf8]{inputenc}
\usepackage[T1]{fontenc}
\usepackage[ngerman]{babel}
\usepackage{graphicx}
\usepackage{array}
\usepackage{booktabs}
\usepackage{float}

\usepackage[scaled]{helvet}
\renewcommand{\familydefault}{\sfdefault}

\usepackage[a4paper, top=2cm, bottom=2cm, left=2cm, right=2cm]{geometry}

\usepackage{setspace}
\setstretch{1.5}

\setlength{\parindent}{0pt}
\setlength{\parskip}{6pt}

\usepackage{microtype}
\sloppy
\hyphenpenalty=1000
\tolerance=3000

\renewcommand{\footnotesize}{\fontsize{10}{12}\selectfont}

\setcounter{secnumdepth}{3}
\setcounter{tocdepth}{3}

\usepackage{titlesec}
\titleformat{\section}{\normalfont\fontsize{12}{14}\bfseries}{\thesection}{1em}{}
\titleformat{\subsection}{\normalfont\fontsize{12}{14}\bfseries}{\thesubsection}{1em}{}
\titleformat{\subsubsection}{\normalfont\fontsize{12}{14}\bfseries}{\thesubsubsection}{1em}{}

\usepackage[
  colorlinks=true,
  linkcolor=black,
  citecolor=blue,
  filecolor=black,
  urlcolor=blue
]{hyperref}
\usepackage[capitalise,nameinlink]{cleveref}

\usepackage{fancyhdr}
\pagestyle{fancy}
\fancyhf{}
\renewcommand{\headrulewidth}{0pt}
\fancyfoot[C]{\thepage}

\usepackage[backend=biber, style=apa]{biblatex}
\addbibresource{references.bib}

\usepackage{titling}

\usepackage{acronym}
\usepackage[german=quotes]{csquotes}

\usepackage{caption}
\captionsetup[table]{
    font=small,
    skip=10pt,
    labelfont=bf
}

\usepackage[table]{xcolor}
\definecolor{zebragray}{gray}{0.92}

\begin{document}

\begin{titlepage}
    \thispagestyle{empty}
    \centering
    \vspace*{5cm}
    {\Huge\bfseries Portfolio \par}
    \vspace{0.3cm}
    {\Large Erarbeitungs-/Reflexionsphase (Phase~2) \par}
    \vspace{1cm}
    {\large Sondierung produktionsreifer KI-Use-Cases für lokale KMU \\ und Skizze eines möglichen Beratungsangebots \par}
    \vspace{1.5cm}
    {\large Kurs: DLBPKIEKPT01 -- Project: AI Fluency -- Einführung in die Generative KI \par}
    \vspace{0.5cm}
    {\large Studiengang: Angewandte Künstliche Intelligenz \par}
    \vspace{0.5cm}
    {\large Sven Behrens \par}
    \vspace{0.5cm}
    {\large Matrikelnummer: 42303511 \par}
    \vspace{0.5cm}
    {\large Tutor: Prof. Dr. Thorsten Fröhlich \par}
    \vspace{0.5cm}
    {\large \today \par}
\end{titlepage}

\pagenumbering{Roman}
\setcounter{page}{1}

\tableofcontents
\newpage

\pagenumbering{arabic}
\setcounter{page}{1}

% =============================================================================
% PHASE 2 -- ERARBEITUNGS-/REFLEXIONSPHASE
% Laut Aufgabenstellung besteht Phase 2 aus zwei Teilen:
%   1. Zwischenprodukt              (keine Seitenvorgabe)
%   2. Dokumentation Arbeitsweise   (5 Beispiele, je 0,5-1 Seite)
%      - 3x Strategische KI-Nutzung (Beispiel A)
%      - 1x Bewusst ohne KI         (Beispiel B)
%      - 1x KI-Fehler gefunden      (Beispiel C)
% =============================================================================

\section{Zwischenprodukt}
% Inhalt: Business Model Canvas + Branchen-Use-Case-Mapping.
% Darf unfertig sein, wird in Phase 3 weiterentwickelt.

[Einleitende Bemerkung folgt: Was wird im Zwischenprodukt geliefert,
welcher Reifegrad wird in dieser Phase angestrebt.]

\subsection{Branchen-Use-Case-Mapping}
% Inhalt: 3-5 Branchen mit jeweils produktionsreifen, eingeschränkt
% einsatzfähigen und ungeeigneten Use-Cases.

[Inhalt folgt.]

\subsection{Business Model Canvas}
% Inhalt: Canvas-Bausteine für das hypothetische KMU-KI-Beratungsangebot.

[Inhalt folgt.]

\section{Dokumentation der Arbeitsweise}

Die folgende Dokumentation gibt fünf konkrete Arbeitssituationen aus der
vorliegenden Sondierung wieder. Drei Beispiele zeigen die strategische
Nutzung generativer KI, ein Beispiel die bewusste Nicht-Nutzung und ein
Beispiel die Identifikation und Korrektur eines KI-Fehlers.

% =============================================================================
% STRUKTUR PRO BEISPIEL (Aufgabenstellung):
%   Beispiel A (Strategische KI-Nutzung):
%     - Die Aufgabe
%     - Der Prompt (mit Begründung)
%     - Der KI-Output (Screenshot oder Textauszug)
%     - Kritische Prüfung (brauchbar / falsch / bias / Vereinfachung)
%     - Verbesserung (was wurde selbst getan)
%   Beispiel B (Bewusst ohne KI):
%     - Die Aufgabe
%     - Warum keine KI
%     - Die Lösung
%   Beispiel C (KI-Fehler gefunden \& korrigiert):
%     - Was war falsch (faktisch / biased / oberflächlich)
%     - Wie erkannt
%     - Wie korrigiert
% =============================================================================

\subsection{Beispiel A1: [Titel folgt]}
\paragraph{Die Aufgabe.} [Inhalt folgt.]
\paragraph{Der Prompt.} [Inhalt folgt.]
\paragraph{Der KI-Output.} [Inhalt folgt.]
\paragraph{Kritische Prüfung.} [Inhalt folgt.]
\paragraph{Verbesserung.} [Inhalt folgt.]

\subsection{Beispiel A2: [Titel folgt]}
\paragraph{Die Aufgabe.} [Inhalt folgt.]
\paragraph{Der Prompt.} [Inhalt folgt.]
\paragraph{Der KI-Output.} [Inhalt folgt.]
\paragraph{Kritische Prüfung.} [Inhalt folgt.]
\paragraph{Verbesserung.} [Inhalt folgt.]

\subsection{Beispiel A3: [Titel folgt]}
\paragraph{Die Aufgabe.} [Inhalt folgt.]
\paragraph{Der Prompt.} [Inhalt folgt.]
\paragraph{Der KI-Output.} [Inhalt folgt.]
\paragraph{Kritische Prüfung.} [Inhalt folgt.]
\paragraph{Verbesserung.} [Inhalt folgt.]

\subsection{Beispiel B: [Titel folgt -- bewusst ohne KI]}
\paragraph{Die Aufgabe.} [Inhalt folgt.]
\paragraph{Warum keine KI.} [Inhalt folgt.]
\paragraph{Die Lösung.} [Inhalt folgt.]

\subsection{Beispiel C: [Titel folgt -- KI-Fehler gefunden \& korrigiert]}
\paragraph{Was war falsch.} [Inhalt folgt.]
\paragraph{Wie erkannt.} [Inhalt folgt.]
\paragraph{Wie korrigiert.} [Inhalt folgt.]

\newpage

\printbibliography
\addcontentsline{toc}{section}{Literaturverzeichnis}

\end{document}
